\documentclass[12pt]{article}
\usepackage[utf8]{inputenc}
\usepackage[T1]{fontenc}
\usepackage{hyperref}
\title{Upper Taromi Language Resources for K-12 Curriculum Development}
\author{Compiled by Claude Sonnet 5 (claude-sonnet-5) for Andr\'as Kornai}
\date{}
\begin{document}
\maketitle

\section{Introduction}
Upper Taromi (ISO 639-3: tov; Glottolog: uppe1235) is a small Northwestern Iranian variety of Iran. Glottolog places it in Tatic (Taromic), while Ethnologue and Verbix place it with Talysh. Joshua Project gives a population of about 300, and Verbix rates the language dormant. Very little exists.

\section{Primary Educational Resources}
None found.

\section{Linguistic and Reference Materials}
\begin{itemize}
\item Verbix endangered languages page: \url{https://docs.verbix.com/EndangeredLanguages/UpperTaromi}
\item ASJP word list: \url{https://asjp.clld.org/languages/UPPER_TAROMI}
\item Wikipedia, Tati language (Iran), which includes Taromi: \url{https://www.Wikipedia.org/wiki/Kabatei}
\end{itemize}

\section{Religious and Community Resources}
\begin{itemize}
\item Joshua Project profile: \url{https://joshuaproject.net/people_groups/18866}
\end{itemize}

\section{Digital and Audio Resources}
Joshua Project notes that audio recordings exist, but I did not find or verify any.

\section{Curriculum Development Context}
No school programme was found. With about 300 speakers and a possibly dormant status, any curriculum would be a documentation and revitalisation effort. Persian is the language of schooling in Iran.
\end{document}